\newcommand\Pal{\mathit{Pal}}
\newcommand\Zag{\mathit{Zag}}
\newcommand\BinSab{\mathit{BinSab}}
\newcommand\Sab{\mathit{Sab}}
\newcommand\Mnoz{\mathit{Mnoz}}
\newcommand\Sloz{\mathit{Sloz}}
\newcommand\Prost{\mathit{Prost}}
\newcommand\Th{\mathrm{Th}}

\chapter{Formalne teorije}

Formalne teorije predstavljaju apstrakciju deduktivnog procesa. U ovom poglavlju
se sre\'cemo sa osnovnim pojmovima formalnih teorija i razmatramo nekoliko
veoma jednostavnih, ve\v sta\v ckih primera kako bismo se navikli na terminologiju
i strategije koje su uobi\v cajene u rezonovanju o formalnim matemati\v ckim
teorijama. U naredna dva poglavlja \'cemo predstaviti dve va\v zne formalne teorije
u matematici: iskazni ra\v cun i predikatski ra\v cun.

\section{Pojam formalne teorije}

\emph{Alfabet}%
    \index{alfabet}
je proizvoljan neprazan kona\v can skup. Elemente alfabeta $A$ \'cemo
zvati \emph{slova}%
    \index{slovo},
a \emph{re\v c nad $A$}%
    \index{re\v c}
je svaki niz simbola iz~$A$.
Preciznije, \emph{re\v c du\v zine $k$ nad alfabetom $A$} je proizvoljan element skupa~$A^k$.
Prilikom zapisivanja re\v ci ne\'cemo pisati zareze i zagrade kako bi re\v ci vi\v se li\v cile na
re\v ci na koje smo navikli.

\paragraph{Primer.}
  Evo nekoliko re\v ci nad alfabetom $A = \{a, b, n\}$: \textit{banana}, \textit{abba},
  \textit{aa}, ili prosto \textit{n}. Prva od njih ima \v sest slova, slede re\v c sa \v cetiri
  slova, re\v c sa dva slova i na kraju re\v c od jednog slova.

\bigskip

Va\v zno je razumeti da su re\v ci sa kojima \'cemo raditi u ovom kursu \emph{formalne re\v ci},
dakle, nizovi simbola kojima nije pridru\v zeno nikakvo zna\v cenje. Prema tome,
\textit{nbbaaa} je jednako dobra re\v c kao \textit{banana}. Problem zna\v cenja
re\v ci \'cemo prepusiti drugim disciplinama.

Neka je~$\lambda$ re\v c nad alfabetom~$A$.
Du\v zinu re\v ci~$\lambda$%
    \index{du\v zina re\v ci}%
    \index{re\v c!du\v zina}
\'cemo ozna\v citi sa~$|\lambda|$%
    \glossary{{$|\lambda|$}{du\v zina re\v ci $\lambda$}}.
Za slovo~$a \in A$ sa $|\lambda|_a$%
    \glossary{{$|\lambda|_a$}{broj pojavljivanja slova $a$ u re\v ci $\lambda$}}
ozna\v cavamo broj pojavljivanja slova~$a$ u re\v ci~$\lambda$.

\paragraph{Primer.}
  Neka je $A = \{a, b, c, n\}$ i neka je $\lambda = banana$ re\v c nad~$A$. Tada je $|\lambda| = 6$,
  $|\lambda|_a = 3$, $|\lambda|_b = 1$, $|\lambda|_c = 0$ i $|\lambda|_n = 2$.

\bigskip

Najjednostavnija i najprirodnija operacija nad formalnim re\v cima je nadovezivanje re\v ci.
Neka su $\lambda$ i $\mu$ dve re\v ci nad alfabetom~$A$. Tada sa $\lambda\mu$ ozna\v cavamo re\v c koja se dobija
tako \v sto re\v ci $\lambda$ i $\mu$ nadove\v zu, odnosno, tako \v sto se slova re\v ci $\mu$
dopi\v su odmah iza slova re\v ci~$\lambda$. Za re\v c $\lambda\mu$ ka\v zemo da je dobijena
\emph{konkatenacijom}%
    \index{konkatenacija}
ili \emph{nadovezivanjem}%
    \index{nadovezivanje re\v ci}
re\v ci~$\lambda$ i~$\mu$.

\paragraph{Primer.} Neka je $A = \{a, b, c, n\}$ i neka je $\lambda = ban$ i $\mu = ana$.
  Tada je $\lambda\mu = banana$.

\bigskip

Dopu\v stamo i re\v ci bez slova. Na svakom alfabetu postoji ta\v cno jedna takva re\v c
koju zovemo \emph{prazna re\v c}%
    \index{prazna re\v c}%
    \index{re\v c!prazna}
i ozna\v cavamo je sa~$\epsilon$%
    \glossary{{$\epsilon$}{prazna re\v c}}.
To je re\v c du\v zine~$0$.
Napomenimo jo\v s da je $\epsilon \lambda = \lambda \epsilon = \lambda$ za svaku re\v c~$\lambda$.
Drugim re\v cima, prazna re\v c je neutralni element za operaciju nadovezivanja re\v ci.
Skup \emph{svih} re\v ci (uklju\v cuju\'ci i praznu re\v c) nad alfabetom $A$ ozna\v cavamo sa $A^*$,%
    \glossary{{$A^*$}{skup svih re\v ci nad alfabetom $A$, uklju\v cuju\'i i praznu re\v c}}
dok skup svih \emph{nepraznih} re\v ci nad alfabetom $A$ ozna\v cavamo sa~$A^+$.%
    \glossary{{$A^+$}{skup svih nepraznih re\v ci nad alfabetom $A$}}
    
\paragraph{Definicija.}
  \emph{Formalna teorija} je ure\dj ena \v cetvorka $T = (A, F, X, R)$ gde je $A$ neki alfabet, $F$ je podskup skupa
  $A^*$ koga zovemo \emph{skup formula}, $X$ je podskup skupa $F$ koga zovemo \emph{skup aksioma},
  a $R$ je skup relacija na $F$ raznih arnosti koje zovemo \emph{pravila izvo\dj enja}.

\paragraph{Primer: Formalna teorija $\Pal$.} Neka je $\Pal = (A, F, X, R)$ formalna teorija \v ciji alfabet je
$A = \{a, d, i, m\}$, gde je $F = A^*$ (dakle, formule su sve re\v ci nad alfabetom~$A$),
aksiome su sve jednoslovne re\v ci $X = \{a, d, i, m\}$, i koja ima \v cetiri pravila izvo\dj enja $R = \{r_a, r_d, r_i, r_m\}$, gde je:
$$
  \begin{array}{l@{\qquad}l}
  r_a = \{(\phi, a\phi a) : \phi \in F\}, & r_i = \{(\phi, i\phi i) : \phi \in F\},\\
  r_d = \{(\phi, d\phi d) : \phi \in F\}, & r_m = \{(\phi, m\phi m) : \phi \in F\}.
  \end{array}
$$
Pravila izvo\dj enja do\v zivljavamo na slede\'ci na\v cin: pravilo $r_a$ zna\v ci da iz formule $\phi$
``mo\v zemo da izvedemo'' formulu $a \phi a$, i sli\v cno za ostala dva pravila. Zato pravila izvo\dj enja
\v ce\v s\'ce (i preglednije) zapisujemo i ovako:
$$
  r_a : \begin{array}{c}
    \phi \\ \hline
    a \phi a
  \end{array}, \qquad
  r_d : \begin{array}{c}
    \phi \\ \hline
    d \phi d
  \end{array}, \qquad
  r_i : \begin{array}{c}
    \phi \\ \hline
    i \phi i
  \end{array}, \qquad
  r_m : \begin{array}{c}
    \phi \\ \hline
    m \phi m
  \end{array}.
$$

\paragraph{Primer: Formalna teorija $\Zag$.} Neka je $\Zag = (A, F, X, R)$ formalna teorija \v ciji alfabet je
$A = \{(, )\}$, gde je
$
  F = \{ \lambda \in A^* : |\lambda|_( = |\lambda|_)\}
$
skup svih re\v ci u kojima je broj otvorenih zagrada jednak broju zatvorenih zagrada,
koja ima samo jednu aksiomu $X = \{()\}$, i dva pravila izvo\dj enja $R = \{r, s\}$, gde je:
$$
  r = \{(\phi, (\phi)) : \phi \in F\}, \qquad s = \{(\phi, \psi, \phi\psi) : \phi, \psi \in F\},
$$
ili, preglednije,
$$
  r : \begin{array}{c}
    \phi \\ \hline
    (\phi)
  \end{array}, \qquad
  s : \begin{array}{c}
    \phi \qquad \psi \\ \hline
    \phi\psi
  \end{array}.
$$
Pravilo $s$ zna\v ci da iz formule $\phi$ i formule $\psi$ ``mo\v zemo da izvedemo'' formulu $\phi\psi$.

\bigskip

Pravila izvo\dj enja u formalnim teorijama nam omogu\'cuju da, polaze\'ci od aksioma formalne teorije, izvedemo neke
od formula formalne teorije (te formule koje se mogu ``izvesti''
smatramo ``dobrim'' ili ``va\v znim'').

%\paragraph{Primer: Formalna teorija $\Pal$ (nastavak).} Evo primera ``izvo\dj enja'' jedne formule u formalnoj
%teoriji $\Pal$:
%$$
%  i \;\overset{r_m}{\longrightarrow}\; mim
%    \;\overset{r_a}{\longrightarrow}\; amima
%    \;\overset{r_d}{\longrightarrow}\; damimad
%    \;\overset{r_a}{\longrightarrow}\; adamimada
%    \;\overset{r_m}{\longrightarrow}\; madamimadam.
%$$

\bigskip

\begin{explanationpic}
  \paragraph{Primer: Formalna teorija $\Pal$ (nastavak).} Pored je dat primer ``izvo\dj enja'' jedne formule u formalnoj
  teoriji $\Pal$.
  
  Pitanje: Da li mo\v zemo da pretpostavimo \v sta daje svako ``izvo\dj enje'' u ovoj formalnoj teoriji?
  (Uputstvo: ``Madam, I'm Adam.'')
\sidepic
  \ \\
  \input ft01-pal.pgf
\end{explanationpic}

\paragraph{Primer: Formalna teorija $\Zag$ (nastavak).} Evo primera ``izvo\dj enja'' jedne formule u formalnoj
  teoriji $\Zag$.

\begin{center}
  \input ft01-zag.pgf
\end{center}

  Pitanje: Da li mo\v zemo da pretpostavimo \v sta daje svako ``izvo\dj enje'' u ovoj formalnoj teoriji?
  (Uputstvo: setiti se algebarskih izraza sa zagradama.)


\bigskip

Kada bi formule u formalnoj teoriji bile neke matemati\v cke tvrdnje onda bi
``izvo\dj enje'' formule zapravo odgovaralo dokazu tvrdnje. Zato se pojam ``izvo\dj enja''
formalizuje na slede\'ci na\v cin.

\paragraph{Definicija.} \emph{Dokaz} u formalnoj teoriji $T$ je niz formula
$
  \phi_1, \, \phi_2, \, \ldots, \, \phi_n
$
teorije $T$ takav da je svaka formula u nizu ili aksioma teorije $T$, ili je dobijena od nekih od
prethodnih formula u nizu primenom nekog od pravila izvo\dj enja teorije~$T$.
Za poslednju formulu u dokazu ka\v zemo da je \emph{posledica} formalne teorije~$T$.
Skup svih posledica formalne teorije~$T$ ozna\v cavamo sa~$\Th(T)$.

\paragraph{Primer: Formalna teorija $\Pal$ (nastavak).}
Poka\v zimo da je $\Th(\Pal)$ skup svih palindroma neparne du\v zine
koji se mogu napisati na alfabetu~$A$.
(Re\v c je \emph{palindrom} ako se i sleva na desno i zdesna na levo \v cita na isti na\v cin;
na primer \textit{abba} i \textit{potop} su palindromi, prvi parne, a drugi neparne du\v zine.)
Dakle, dokazujemo slede\'cu jednakost skupova:
$$
  \Th(\Pal) = \{ \lambda \in A^* : \lambda \text{ je palindrom neparne du\v zine}\}.
$$

$(\subseteq)$ Neka je $\lambda \in \Th(\Pal)$. Tada postoji dokaz u kome je $\lambda$
poslednja formula. Dokazujemo da je $\lambda$ palindrom neparne du\v zine
indukcijom po du\v zini dokaza za~$\lambda$.

Ako postoji dokaz za $\lambda$ du\v zine~1, ovako:
$$
  \lambda
$$
onda je $\lambda$ aksioma (jer ne mo\v ze biti izvedena iz nekih prethodnih formula primenom nekog od
pravila izvo\dj enja), a sve aksiome teorije $\Pal$ su palnidromi neparne du\v zine.

Pretpostavimo da je tvr\dj enje ta\v cno za sve posledice teorije $\Pal$ koje imaju dokaz du\v zine
manje od $n$ i neka je $\lambda$ posledica teorije $\Pal$ koja ima dokaz du\v zine~$n$:
$$
  \phi_1, \, \phi_2, \, \ldots, \, \phi_{n-1}, \, \lambda.
$$
Ako je $\lambda$ aksioma, tvr\dj enje sledi odmah. Pretpostavimo, zato, da $\lambda$ nije aksioma.
Onda postoji neko $k < n$ i neko pravilo izvo\dj enja $r_x$ takvo da je formula $\lambda$ dobijena
od formule $\phi_k$ primenom pravila $r_x$. Odatle sledi da je
$$
  \lambda = x \phi_k x.
$$
Obzirom da formula $\phi_k$ ima dokaz du\v zine $k < n$, ona je palindrom neparne du\v zine
prema induktivnoj hipotezi, pa je onda
i formula $x \phi_k x = \lambda$ palindrom neparne du\v zine.

\medskip

$(\supseteq)$ Neka je $\lambda \in A^*$ palindrom neparne du\v zine. Poka\v zimo indukcijom po du\v zini
re\v ci $\lambda$ da je $\lambda \in \Th(\Pal)$.

Ako je $|\lambda| = 1$ onda je $\lambda$ re\v c du\v zine 1 nad alfabetom $A$, a sve takve re\v ci su
aksiome teorije $\Pal$. Zato je $\lambda \in \Th(\Pal)$.

Pretpostavimo da je tvr\dj enje ta\v cno za sve palindrome neparne du\v zine koja je strogo
manja od $n$ i neka je $\lambda$ palindrom neparne du\v zine~$n$. Kako je $\lambda$ palindrom,
postoji slovo $x \in A$ i re\v c $\mu \in A^*$ takvi da je $\lambda = x \mu x$. Zato \v sto je $\lambda$
palindrom neparne du\v zine lako se zaklju\v cuje da je $\mu$ tako\dj e palindrom i da je
$|\mu| = |\lambda| - 2$ neparan broj manji od $n$. Prema induktivnoj hipotezi, postoji dokaz
u $T$ \v cija poslednja formula je $\mu$:
$$
  \phi_1, \; \phi_2, \; \ldots, \; \phi_l, \; \mu.
$$
Sada se lako vidi da je
$$
  \phi_1, \; \phi_2, \; \ldots, \; \phi_l, \; \mu, \; \underbrace{x \mu x}_{= \lambda}
$$
tako\dj e dokaz u $T$ jer se poslednja formula u ovom dokazu dobija od pretposlednje primenom pravila $r_x$.

Ovim je okon\v can dokaz tvr\dj enja da je $\Th(\Pal)$ skup svih palindroma neparne du\v zine
koji se mogu napisati na alfabetu~$A$.

\paragraph{Primer: Formalna teorija $\Zag$ (nastavak).}
Poka\v zimo da je $\Th(\Zag)$ skup svih \emph{pravilnih zagradastih re\v ci}.

\emph{Zagradasta re\v c} je svaka re\v c nad alfabetom $A = \{ (, ) \}$. Za zagradastu re\v c $\lambda$
ka\v zemo da je \emph{pravilna} ako
\begin{itemize}
\item
  ima isti broj otvorenih i zatvorenih zagrada, i
\item
  u svakom po\v cetnom segmentu je broj otvorenih zagrada $\ge$ od broja zatvorenih zagrada.
\end{itemize}
Na primer, $))())()$ je jedna zagradasta re\v c. Ona nije pravilna jer nema isti broj otvorenih i zatvorenih zagrada.
Zagradasta re\v c $())(()$ ima isti broj otvorenih i zatvorenih zagrada, ali nije pravilna jer po\v cetni segment
$())$ ne ispunjava drugi zahtev. Primer pravilne zagradaste re\v ci bi bio, recimo, $(())()(()())$. Intuitivno,
u pravilnoj zagradastoj re\v ci su zagrade raspore\dj ene ``pravilno'' u smislu raspore\dj ivanja zagrada u
algebarskom izrazu. Svaka pravilna zagradasta re\v c se mo\v ze dobiti od nekog pravilno zapisanog algebarskog izraza
kada iz njega izbacimo sve simbole koji nisu zagrade.
Dakle, poka\v zimo da je:
$$
  \Th(\Zag) = \{ \lambda \in F : \lambda \text{ je pravilna zagradasta re\v c} \}.
$$

$(\subseteq)$ Neka je $\lambda \in \Th(\Zag)$. Tada postoji dokaz u kome je $\lambda$
poslednja formula. Dokazujemo da je $\lambda$ pravilna zagradasta re\v c
indukcijom po du\v zini dokaza za~$\lambda$.

Ako postoji dokaz za $\lambda$ du\v zine~1
onda je $\lambda$ aksioma (jer ne mo\v ze biti izvedena iz nekih prethodnih formula primenom nekog od
pravila izvo\dj enja), a jedina aksiome teorije $\Zag$ je $()$, \v sto je pravilna zagradasta re\v c.

Pretpostavimo da je tvr\dj enje ta\v cno za sve posledice teorije $\Zag$ koje imaju dokaz du\v zine
manje od $n$ i neka je $\lambda$ posledica teorije $\Zag$ koja ima dokaz du\v zine~$n$:
$$
  \phi_1, \, \phi_2, \, \ldots, \, \phi_{n-1}, \lambda.
$$
Ako je $\lambda$ aksioma, tvr\dj enje sledi odmah. Pretpostavimo, zato, da $\lambda$ nije aksioma.
Onda je $\lambda$ dobijena primenom nekog od pravila izvo\dj enja na neke od formula koje se javljaju
ranije u dokazu. Obzirom da teorija ima dva pravila izvo\dj enja, razmatramo dva slu\v caja:

\medskip

$(1^\circ)$ $\lambda$ je dobijena primenom pravila $r$ na formulu $\phi_j$ za neko $j < n$. Tada je
$\lambda = (\phi_j)$. Formula $\phi_j$ je pravilna zagradasta re\v c prema induktivnoj hipotezi, pa odatle
odmah sledi da je i $\lambda = (\phi_j)$ pravilna zagradasta re\v c.

\medskip

$(2^\circ)$ $\lambda$ je dobijena primenom pravila $s$ na formule $\phi_j$  i $\phi_k$
za neke $j, k < n$. Tada je $\lambda = \phi_j \phi_k$.
Formule $\phi_j$ i $\phi_k$ su pravilne zagradaste re\v ci prema induktivnoj hipotezi, pa odatle
odmah sledi da je i $\lambda = \phi_j \phi_k$ pravilna zagradasta re\v c.

\medskip

$(\supseteq)$ Neka je $\lambda \in F$ pravilna zagradasta re\v c. Poka\v zimo indukcijom po du\v zini
re\v ci $\lambda$ da je $\lambda \in \Th(\Zag)$.

Nijedna re\v c du\v zine 1 ne mo\v ze biti pravilna zagradasta re\v c, tako da indukciju po\v cinjemo
od $|\lambda| = 2$. Jedina pravilna zagradasta re\v c du\v zine~2 je $()$, \v sto je
aksioma teorije $\Zag$. Zato je $\lambda \in \Th(\Zag)$.

Pretpostavimo da je tvr\dj enje ta\v cno za sve pravilne zagradaste re\v ci \v cina du\v zina je strogo
manja od $n$ i neka je $\lambda$ pravilna zagradasta re\v c du\v zine~$n$. Tada, prema definiciji
pravilne zagradaste re\v ci, prvo slovo re\v ci $\lambda$ mora biti~$($. Uo\v cimo zatvorenu zagradu
koja odgovara ovoj otvorenoj zagradi. Ovim parom odgovaraju\'cih zagrada je re\v c $\lambda$ razlo\v zena
na slede\'ci na\v cin:
$$
  \lambda = (\mu)\sigma
$$
Zato \v sto je $\lambda$ pravilna zagradasta re\v ci i zato \v sto je uo\v ceni par zagrada par
odgovaraju\'cih zagrada sledi da svaka od re\v ci $\mu$ i $\sigma$ prazna re\v c ili pravilna zagradasta
re\v c du\v zine manje od $|\lambda|$. Sada bi trebalo razmotriti \v cetiri slu\v caja u zavisnosti od
toga da li su $\mu$ i $\sigma$ prazne re\v ci. Mi razmotriti samo najinteresantniji slu\v caj
kada ni $\mu$ ni $\sigma$ nisu prazne re\v ci jer se dokazi u preostala tri slu\v caja mogu dobiti
modifikacijom ovog dokaza.

Prema induktivnoj hipotezi postoje dokazi u $T$ \v cija poslednje formule su $\mu$, odnosno,~$\sigma$:
$$
  \phi_1, \; \phi_2, \; \ldots, \; \phi_k, \; \mu \text{\qquad i\qquad}
  \psi_1, \; \psi_2, \; \ldots, \; \psi_l, \; \sigma.
$$
Nadovezivanjem ova dva dokaza se lako dobija dokaz formule $\lambda$:
$$
  \phi_1, \; \phi_2, \; \ldots, \; \phi_k, \; \mu, \;
  \psi_1, \; \psi_2, \; \ldots, \; \psi_l, \; \sigma, \;
  (\mu), \; \underbrace{(\mu)\sigma}_{= \lambda},
$$
gde smo pretposlednju formulu novog dokaza dobili tako \v sto smo na formulu $\mu$ primenili pravilo $r$,
a poslednju tako \v sto smo na pretposlednju formule $(\mu)$ i $\sigma$
primenili pravilo~$s$.

Ovim je okon\v can dokaz tvr\dj enja da je $\Th(\Zag)$ skup svih pravilnih zagradastih re\v ci.

\paragraph{Primer.} Na\'ci formalnu teoriju nad alfabetom $A = \{0, 1\}$ \v cije posledice su sve re\v ci
kod kojih je broj jedinica paran.

\medskip

\noindent
\textit{Re\v senje.} To je formalna teorija $T = (A, F, X, R)$ gde je $A = \{0, 1\}$, $F = A^*$,
$X = \{\epsilon\}$ i koja ima dva pravila izvo\dj enja:
$$
  r : \begin{array}{c}
    \phi \\ \hline
    11\phi
  \end{array}, \qquad
  s : \begin{array}{c}
    \phi\psi \\ \hline
    \phi0\psi
  \end{array}.
$$
Pravilo $s$ zna\v ci da se iz formule $\lambda = \phi\psi$ mo\v ze izvesti formula $\mu = \phi0\psi$,
odnosno, da se u formuli $\lambda$ mo\v ze dopisati nula gdegod to po\v zelimo.

\paragraph{Primer: Sve mo\v ze da se opi\v se formalnom teorijom.}
Za dati proizvoljan alfabet $A$ i proizvoljno $S \subseteq A^*$ na\'ci formalnu teoriju $T$ nad alfabetom
$A$ takvu da je $\Th(T) = S$.

\medskip

\noindent
\textit{Re\v senje.} To je formalna teorija $T = (A, F, X, R)$ gde je $A$ dato, $F = A^*$,
$X = S$ i koja ima samo jedno pravilo izvo\dj enja:
$$
  r : \begin{array}{c}
    \phi \\ \hline
    \phi
  \end{array}.
$$
Ova formalna teorija ne mo\v ze ni\v sta novo da ``izvede'' jer su svi dokazi oblika
$$
  \phi, \; \phi, \; \ldots, \; \phi.
$$
Zato je $\Th(T) = X = S$.

\bigskip

Prethodni primer pokazuje za\v sto je va\v zno insistirati na tome da skupovi aksioma ``va\v znih'' teorija
imaju neka specijalna svojstva. Da bi teorija bila ``interesantna'' i ``korisna'' va\v zno je da njen skup aksioma
bude ``minimalan'', ``efektivno nabrojiv'' ili na neki drugi na\v cin ``lep''.

\section{Formalne teorije umeju da ra\v cunaju}

Postoje situacije u kojima je proces izvo\dj enja formule interesantniji od opisa skupa svih formula
koje su posledice formalne teorije. Primer koga pokazujemo u ovom odeljku ilustruje nekoliko takvih situacija.
Pokaza\'cemo formalnu teoriju koja sabira brojeve, formalnu teoriju koja mno\v zi brojeve,
formalnu teoriju koja opisuje sve slo\v zene brojeve i
formalnu teoriju koja opisuje sve proste brojeve. Sve formalne teorije koje ovde pokazujemo
je predlo\v zio Daglas Hof\v stater u svojoj knjizi ``Gedel, E\v ser, Bah: Beskrajna zlatna pletenica''\footnote{%
  Douglas R. Hofstadter: ``G\"odel, Escher, Bach: An Eternal Golden Braid'',
  Basic Books 1999 (first published 1979)}. Mi smo samo prilagodili oznake i prezentaciju ovih teorija.

\begin{quote}
  \textbf{Ovo je samo prvi, mali korak ka osnovnom cilju kursa:
  demonstrirati da formalne teorije mogu da izvode proizvoljne ra\v cunske procese.}
\end{quote}


\subsection{Formalna teorija $\Sab$}

Formalna teorija $\Sab$ opisuje proces sabiranja prirodnih brojeva. Brojeve \'cemo predstavljati kao nizove recki, ovako:
\begin{align*}
  &1 \quad\to\quad !\\
  &2 \quad\to\quad !!\\
  &3 \quad\to\quad !!!
\end{align*}
Da bismo napravili jasnu razliku izme\dj u simbola formalne teorije i simbola metajezika, aritmeti\v cke izraze \'cemo zapisivati ovako:
\begin{align*}
  &2 + 3 = 5 \quad\to\quad !! \oplus !!! \equiv !!!!!
\end{align*}
Neka je $\Sab = (A, F, X, R)$ formalna teorija nad alfabetom $A = \{!, \Boxed\oplus, \Boxed\equiv\}$ \v ciji skup formula
je $F = A^*$, \v ciji skup aksioma je
$$
  X = \{\phi \oplus ! \equiv \phi! : \phi \in \{!\}^+\}
$$
i koja ima jedno pravilo izvo\dj enja:
$$
  \begin{array}{c}
    \phi \oplus \psi \equiv \theta \\ \hline
    \phi \oplus \psi! \equiv \theta!
  \end{array}, \text{\quad za sve } \phi, \psi, \theta \in \{!\}^+.
$$
Tada se lako vidi da je
$$
  \Th(\Sab) = \{\phi \oplus \psi \equiv \theta : \phi, \psi, \theta \in \{!\}^+ \text{ i } |\phi| + |\psi| = |\theta|\}.
$$
Dokaz se zasniva na \v cinjenici da je $x + (y + 1) = (x + y) + 1$.

\subsection{Formalna teorija $\Mnoz$}

Sada \'cemo opisati formalnu teoriju $\Mnoz$ koja opisuje proces mno\v zenja prirodnih brojeva.
Brojeve \'cemo, kao i ranije, predstavljati kao nizove recki, a odgovaraju\'ce aritmeti\v cke izraze \'cemo zapisivati ovako:
\begin{align*}
  &2 \cdot 3 = 6 \quad\to\quad !! \otimes !!! \equiv !!!!!!
\end{align*}
Neka je $\Mnoz = (A, F, X, R)$ formalna teorija nad alfabetom $A = \{!, \Boxed\otimes, \Boxed\equiv\}$ \v ciji skup formula je
$F = A^*$, \v ciji skup aksioma je
$$
  X = \{\phi \otimes ! \equiv \phi : \phi \in \{!\}^+\}
$$
i koja ima jedno pravilo izvo\dj enja:
$$
  \begin{array}{c}
    \phi \otimes \psi \equiv \theta \\ \hline
    \phi \otimes \psi! \equiv \theta\phi
  \end{array}, \text{\quad za sve } \phi, \psi, \theta \in \{!\}^+.
$$
Opet se lako vidi da je
$$
  \Th(\Mnoz) = \{\phi \otimes \psi \equiv \theta : \phi, \psi, \theta \in \{!\}^+ \text{ i } |\phi| \cdot |\psi| = |\theta|\}.
$$
Dokaz se zasniva na \v cinjenici da je $x \cdot (y + 1) = xy + x$.

\subsection{Formalna teorija $\Sloz$}

Formalna teorija $\Sloz$ je pro\v sirenje formalne teorije $\Mnoz$ koja nam omogu\'cuje da opi\v semo \v cinjenicu da je neki
broj slo\v zen. Da se podsetimo, broj $n$ je slo\v zen ako postoje $a, b \ge 2$ takvi da je $n = ab$.

Neka je $\Sloz = (A, F, X, R)$ formalna teorija nad alfabetom $A = \{!, \Boxed\otimes, \Boxed\equiv, S\}$ \v ciji skup formula je
$F = A^*$, \v ciji skup aksioma je
$$
  X = \{\phi \otimes ! \equiv \phi : \phi \in \{!\}^+\}
$$
i koja ima dva pravila izvo\dj enja:
$$
  \begin{array}{c}
    \phi \otimes \psi \equiv \theta \\ \hline
    \phi \otimes \psi! \equiv \theta\phi
  \end{array}, \text{\quad za sve } \phi, \psi, \theta \in \{!\}^+.
$$
i
$$
  \begin{array}{c}
    \phi! \otimes \psi! \equiv \theta \\ \hline
    S\theta
  \end{array}, \text{\quad za sve } \phi, \psi, \theta \in \{!\}^+.
$$
Sada kada znamo sve posledice teorije $\Mnoz$, lako se vidi da je
\begin{align*}
  \Th(\Sloz) = \mathstrut
    &\{\phi \otimes \psi \equiv \theta : \phi, \psi, \theta \in \{!\}^+ \text{ i } |\phi| \cdot |\psi| = |\theta|\}\\
    &\mathstrut \union \{S\theta : \theta \in \{!\}^+ \text{ i } |\theta| \text{ je slo\v zen broj} \}.
\end{align*}

\subsection{Formalna teorija $\Prost$}

Kona\v cno, pokaza\'cemo i formalnu teoriju koja mo\v ze da opi\v se \v cinjenicu da je broj prost.
Da se podsetimo, broj $n \ge 2$ je prost ako ima samo dva delioca, 1 i $n$. Ako je $n \ge 3$ onda mo\v zemo re\'ci i ovako:
broj $n$ je prost ako nije deljiv nijednim od brojeva 2, 3, \ldots,  $n - 1$.

Neka je $\Prost = (A, F, X, R)$ formalna teorija nad alfabetom $A = \{!, \#, @, P\}$ \v ciji skup formula je
$F = A^*$, \v ciji skup aksioma je
$$
  X = \{\phi \psi \# \phi : \phi, \psi \in \{!\}^+\} \union \{P!!\}
$$
i \v cija pravila izvo\dj enja su:
$$
  r_1 : \begin{array}{c}
    \phi \# \psi \\ \hline
    \phi \# \phi\psi
  \end{array}, \qquad
  r_2 : \begin{array}{c}
    !! \# \phi \\ \hline
    \phi @ !!
  \end{array},\qquad
  r_3 : \begin{array}{c}
    \theta @ \phi \qquad \phi! \# \theta\\ \hline
    \theta @ \phi!
  \end{array},\qquad
  r_4 : \begin{array}{c}
    \phi! @ \phi \\ \hline
    P\phi!
  \end{array},
$$
za sve $\phi, \psi, \theta \in \{!\}^+$.

Pre nego \v sto nastavimo sa analizom ovog sistema da\'cemo intuiciju koju opisuju ova pravila.
Formula
$$
  \phi \# \psi
$$
zna\v ci ``broj $|\phi|$ se ne sadr\v zi u broju $|\psi|$''.
Na primer, aksioma $!!!!\#!!!$ zna\v ci ``4 se ne sadr\v zi u 3'' ili, kra\'ce, $4 \nmid 3$.
Formula
$$
  \phi @ \psi
$$
zna\v ci ``broj $|\phi|$ nije deljiv nijednim od brojeva $2, \ldots, |\psi|$''.
Na primer, formula $!!!!!!!@!!!!!$ zna\v ci ``7 nije deljiv nijednim od brojeva $2, 3, 4, 5$. Kona\v cno, formula
$$
  P\phi
$$
zna\v ci ``$|\phi|$ je prost broj''.

Sada je lako protuma\v citi aksiome: poslednja ka\v ze da je 2 prost broj, dok ostale aksiome zna\v ce da $(n + k) \nmid n$
za sve $n, k \ge 1$.

Pravilo $r_1$ ka\v ze: ``ako $n \nmid k$ onda $n \nmid (n + k)$''.

Pravilo $r_2$ ka\v ze: ``ako $2 \nmid n$ onda $n$ nije deljiv brojem 2''.

Pravilo $r_3$ ka\v ze: ``ako $n$ nije deljiv nijednim od brojeva 2, \ldots, $k$ i ako $(k + 1) \nmid n$ onda
$n$ nije deljiv nijednim od brojeva 2, \ldots, $k + 1$.

Kona\v cno, pravilo $r_4$ ka\v ze: ``ako $n + 1$ nije deljiv nijednim od brojeva 2, \ldots, $n$ onda je $n + 1$ prost broj''.

Sada se mo\v ze pokazati da za $\phi \in \{!\}^+$ va\v zi slede\'ce:
$$
  P\phi \in \Th(\Prost) \text{ ako i samo ako je } |\phi| \text{ prost broj}.
$$
Ovaj dokaz nije te\v zak, ali je tehni\v cki zahtevan. Zato \'cemo umesto dokaza dati primer izvo\dj enja teoreme
$P!!!!!$ koja ka\v ze da je 5 prost broj.

\begin{center}
  \input ft01-prost.pgf
\end{center}

\paragraph{Zadaci.}

\ZAD{ft01-01}{%
  Opisati sve posledice formalne teorije $T = (A, F, X, R)$ \v ciji alfabet je
  $A = \{a, b\}$, gde je $F = A^*$,
  aksioma je samo formula $b$ i koja ima dva pravila izvo\dj enja:
  $$
    r_1 : \begin{array}{c}
      \phi \\ \hline
      a \phi b
    \end{array}, \qquad
    r_2 : \begin{array}{c}
      \phi \\ \hline
      \phi b
    \end{array}.
  $$
}

\ZAD{ft01-02a}{%
  Opisati sve posledice formalne teorije $T = (A, F, X, R)$ \v ciji alfabet je
  $A = \{a, b\}$, gde je $F = A^*$,
  aksioma je samo formula $a$ i koja ima dva pravila izvo\dj enja:
  $$
    r_1 : \begin{array}{c}
      \phi \\ \hline
      \phi a
    \end{array}, \qquad
    r_2 : \begin{array}{c}
      \phi \\ \hline
      \phi b
    \end{array}.
  $$
}

\ZAD{ft01-02b}{%
  Opisati sve posledice formalne teorije $T = (A, F, X, R)$ \v ciji alfabet je
  $A = \{a, b\}$, gde je $F = A^*$,
  aksioma je samo formula $b$ i koja ima dva pravila izvo\dj enja:
  $$
    r_1 : \begin{array}{c}
      \phi \\ \hline
      a \phi
    \end{array}, \qquad
    r_2 : \begin{array}{c}
      \phi \\ \hline
      b \phi
    \end{array}.
  $$
}

\ZAD{ft01-03a}{%
  Opisati sve posledice formalne teorije $T = (A, F, X, R)$ \v ciji alfabet je
  $A = \{a, b, \ldots, z, 0, 1, \ldots, 9\}$ (sva mala slova engleske abecede i sve cifre),
  gde je $F = A^*$, aksiome su $X = \{a, b, \ldots, z\}$ i koja za svako slovo $\sigma \in A$
  ima po jedno pravilo izvo\dj enja:
  $$
    r_\sigma : \begin{array}{c}
      \phi \\ \hline
      \phi \sigma
    \end{array}.
  $$
}

\ZAD{ft01-03b}{%
  Opisati sve posledice formalne teorije $T = (A, F, X, R)$ \v ciji alfabet je
  $A = \{0, 1, \ldots, 9, -, .\}$ (sve cifre, simbol ``minus'' i simbol ``ta\v cka''),
  gde je $F = A^*$, aksiome su $X = \{0, 1, \ldots, 9\}$ i koja za svako $\sigma \in \{0, 1, \ldots, 9\}$
  ima po jedno pravilo izvo\dj enja:
  $$
    r_\sigma : \begin{array}{c}
      \phi \\ \hline
      \phi \sigma
    \end{array},
  $$
  i jo\v s dva pravila izvo\dj enja:
  $$
    r_. : \begin{array}{c}
      \phi \text{\ je formula koja ne sadr\v zi simbol . (ta\v cka)} \\ \hline
      \phi.
    \end{array}
  $$
  $$
    r_- : \begin{array}{c}
      \phi \text{\ je formula koja ne sadr\v zi simbol $-$ (minus)}\\ \hline
      -\phi
    \end{array}.
  $$
}

\ZAD{ft01-04}{%
  Opisati formalnu teoriju nad alfabetom $A = \{0, 1\}$ \v cije posledice su sve re\v ci oblika
  $
    \underbrace{0\ldots 0}_k \underbrace{11\ldots 11}_{2k}
  $
  i ni\v sta drugo, $k \ge 1$.
}

\ZAD{ft01-04a}{%
  Opisati formalnu teoriju nad alfabetom $A = \{0, 1\}$ \v cije posledice su sve re\v ci kod kojih
  se 0 pojavljuje paran broj puta.
}

\ZAD{ft01-04b}{%
  Opisati formalnu teoriju nad alfabetom $A = \{0, 1\}$ \v cije posledice su sve re\v ci kod kojih
  se 0 pojavljuje paran broj puta, a 1 neparan broj puta.
}

\ZAD{ft01-07a}{%
  Opisati \textbf{jednom re\v cenicom} teoreme formalne teorije
  nad alfabetom $\{a, b\}$ \v cije formule su sve re\v ci nad alfabetom, koja ima samo jednu
  aksiomu~$ba$, i koja ima samo jedno pravilo izvo\dj enja:
  $$
  \begin{array}{c}
    \varphi \\ \hline
    ba \varphi
  \end{array}.
  $$
}

\ZAD{ft01-07b}{%
  Opisati \textbf{jednom re\v cenicom} teoreme formalne teorije
  nad alfabetom $\{e, m\}$ \v cije formule su sve re\v ci nad alfabetom, koja ima samo jednu
  aksiomu~$me$, i koja ima samo jedno pravilo izvo\dj enja:
  $$
  \begin{array}{c}
    \varphi \\ \hline
    m \varphi e
  \end{array}.
  $$
}

\ZAD{ft01-07d}{%
  Neka je $(A, F, X, R)$ formalna teorija nad alfabetom $A = \{!, \Boxed\otimes, \Boxed\equiv, D\}$ \v ciji skup formula je
  $F = A^*$, \v ciji skup aksioma je
  $$
    X = \{\phi \otimes ! \equiv \phi : \phi \in \{!\}^+\}
  $$
  i koja ima dva pravila izvo\dj enja:
  $$
    \begin{array}{c}
      \phi \otimes \psi \equiv \theta \\ \hline
      \phi \otimes \psi! \equiv \theta\phi
    \end{array}, \text{\quad za sve } \phi, \psi, \theta \in \{!\}^+.
  $$
  i
  $$
    \begin{array}{c}
      !!!!! \otimes \psi \equiv \theta \\ \hline
      D\theta
    \end{array}, \text{\quad za sve } \psi, \theta \in \{!\}^+.
  $$
  Opisati posledice ove formalne teorije koje su oblika $Q\theta$.
}

\ZAD{ft01-07d}{%
  Neka je $(A, F, X, R)$ formalna teorija nad alfabetom $A = \{!, \Boxed\otimes, \Boxed\equiv, Q\}$ \v ciji skup formula je
  $F = A^*$, \v ciji skup aksioma je
  $$
    X = \{\phi \otimes ! \equiv \phi : \phi \in \{!\}^+\}
  $$
  i koja ima dva pravila izvo\dj enja:
  $$
    \begin{array}{c}
      \phi \otimes \psi \equiv \theta \\ \hline
      \phi \otimes \psi! \equiv \theta\phi
    \end{array}, \text{\quad za sve } \phi, \psi, \theta \in \{!\}^+.
  $$
  i
  $$
    \begin{array}{c}
      \psi \otimes \psi \equiv \theta \\ \hline
      Q\theta
    \end{array}, \text{\quad za sve } \psi, \theta \in \{!\}^+.
  $$
  Opisati posledice ove formalne teorije koje su oblika $Q\theta$.
}


\ZAD{ft01-05a}{%
  Da li je re\v c $ccbaaccbcbabcbaba$ teorema formalne teorije nad alfabetom $\{a, b, c\}$
  \v cije formule su sve re\v ci nad alfabetom, \v ciji skup aksioma je $\{a, b, c\}$ i \v cija pravila izvo\dj enja su
  $$
    \begin{array}{c}
      \varphi \\ \hline
      ba \varphi c
    \end{array}, \quad
    \begin{array}{c}
      \varphi \\ \hline
      ac \varphi b
    \end{array}, \quad
    \begin{array}{c}
      \varphi \\ \hline
      c \varphi ba
    \end{array}?
  $$
}

\ZAD{ft01-05x}{%
  Da li je re\v c $1011001011010101011000110$ posledica formalne teorije nad alfabetom $\{0, 1\}$
  \v cije formule su sve re\v ci nad alfabetom, \v ciji skup aksioma je $\{0, 011, 101\}$ i \v cija pravila izvo\dj enja su
  $$
    \begin{array}{c}
      \varphi \\ \hline
      1\varphi 10
    \end{array}, \quad
    \begin{array}{c}
      \varphi \\ \hline
      101\varphi
    \end{array}, \quad
    \begin{array}{c}
      \varphi \\ \hline
      \varphi011
    \end{array}, \quad
    \begin{array}{c}
      \varphi \qquad \psi \\ \hline
      \varphi 0 \psi
    \end{array}?
  $$
}

\ZAD{ft01-06a}{%
  Na\'ci dokaz posledice $001011010011$ formalne teorije nad alfabetom $\{0, 1\}$ \v cije formule su sve re\v ci,
  koja ima samo jednu aksiomu $01$ i dva pravila izvo\dj enja:
  $$
  \begin{array}{c}
    \varphi \quad \psi \\ \hline
    \varphi\psi
  \end{array}, \quad
  \begin{array}{c}
    \varphi \\ \hline
    0 \varphi 1
  \end{array}.
  $$
}

\ZAD{ft01-MIU}{%
  (Daglas Hof\v stater, 1979)
  Formalna teorija MIU ima alfabet $A = \{M, I, U\}$, formule su sve re\v ci nad $A$, ima samo jednu
  aksiomu, $MI$, i \v cetiri pravila izvo\dj enja:
  $$
  \begin{array}{c}
    \varphi I \\ \hline
    \varphi IU
  \end{array}, \quad
  \begin{array}{c}
    M\varphi \\ \hline
    M \varphi\varphi
  \end{array}, \quad
  \begin{array}{c}
    \varphi III \psi\\ \hline
    \varphi U \psi
  \end{array}, \quad
  \begin{array}{c}
    \varphi UU \psi\\ \hline
    \varphi \psi
  \end{array}.
  $$

  $(a)$ Dokazati da su $MIUIUIUIU$ i $MUIIIIU$ teoreme ovog sistema.

  $(b)$ Da li je $MU$ teorema ovog sistema?
}
% In the MIU formal system introduced by Hofstadter in [Hof79], every word Mx where x is a sequence of U's and I's
% having a number of I's which is not a multiple of 3, is a theorem.


\endinput
%%%%%%%%%%%%%%%%%%%%%%%%%%%%%%%%%%%%%%%%%%%%%%%%%%%%%%%%%%%%%%%%%%%%%%%
%%%%%%%%%%%%%%%%%%%%%%%%%%%%%%%%%%%%%%%%%%%%%%%%%%%%%%%%%%%%%%%%%%%%%%%


\section{$\lambda$-ra\v cun kao formalna teorija}

!!! $\lambda$-ra\v cun 


\endinput
%%%%%%%%%%%%%%%%%%%%%%%%%%%%%%%%%%%%%%%%%%%%%%%%%%%%%%%%%%%%%%%%%%%%%%%
%%%%%%%%%%%%%%%%%%%%%%%%%%%%%%%%%%%%%%%%%%%%%%%%%%%%%%%%%%%%%%%%%%%%%%%

!!!Ovo je stara formalna teorija koja sabira decimalne brojeve cifru po cifru:

Neka je $\Sab = (A, F, X, R)$ formalna teorija nad alfabetom $A = \{\Boxed\oplus, 0, 1, \ldots, 9\}$ \v ciji skup formula
je $F = \{\lambda \in A^* : |\lambda|_\oplus \le 1\}$, \v ciji skup aksioma je $X = \{0, 1, \ldots, 9\}^+$
i koja ima slede\'ca pravila izvo\dj enja.

Pravilo $\mathit{input}$ se mo\v ze primeniti samo na dve aksiome $\alpha, \beta \in X$ iste du\v zine.
Pravilo $\mathit{input}$ preuredi slova ove dve aksiome, postavi inicijalni prenos i marker $\oplus$ na kraj, i tako
omogu\'ci primenu ostalih pravila koja realizuju sabiranje:
$$
  \mathit{input} : \begin{array}{c}
    \alpha \qquad \beta \\ \hline
    \xi 0 \oplus
  \end{array}, \text{\quad samo za } \alpha, \beta \in X \text{ takve da je } |\alpha| = |\beta|,
$$
gde je $\xi$ formula koja za $\alpha = a_1 a_2 \ldots a_k$ i $\beta = b_1 b_2 \ldots b_k$ izgleda ovako:
$$
  \xi = a_1 b_1 a_2 b_2 \ldots a_k b_k.
$$


Pravila $\mathit{output}_0$ i $\mathit{output}_1$ bri\v su marker $\oplus$ sa po\v cetka re\v ci kada se sabiranje zavr\v si.
Ispred markera ostaje prenos poslednjeg zabiranja koji se ignori\v se ako se radi o nuli:
$$
  \mathit{output}_0 : \begin{array}{c}
    0 \oplus \lambda\\ \hline
    \lambda
  \end{array}, \qquad
  \mathit{output}_1 : \begin{array}{c}
    1 \oplus \lambda\\ \hline
    1\lambda
  \end{array},
$$
gde je $\lambda$ neka formula.

Pored ova tri, teorija $\Sab$ ima jo\v s 200 pravila koja (sre\'com!) imaju sli\v cnu strukturu. Za sve
$c, d \in \{0, 1, \ldots, 9\}$ (odgovaraju\'ce cifre prvog i drugog sabirka) i svako
$p \in \{0, 1\}$ (prenos), pravilo $r_{cdp}$ izgleda ovako:
$$
  r_{cdp} : \begin{array}{c}
    \lambda cdp \oplus \mu\\ \hline
    \lambda a \oplus b \mu
  \end{array},
$$
gde su $\lambda$ i $\mu$ neke formule, a $a, b \in \{0, 1, \ldots, 9\}$ su takvi da je $c + d + p = 10 a + b$.
Intuitivno, simbol $\oplus$ sabere tri cifre koje mu prethode, cifru zbira ``gurne'' iza sebe, a cifru prenosa
ostavi ispred sebe kako bi u\v cestvovala u narednom sabiranju. Na primer, pravilo $r_{841}$ izgleda ovako:
$$
  r_{841} : \begin{array}{c}
    \lambda 841 \oplus \mu\\ \hline
    \lambda 1 \oplus 3 \mu
  \end{array},
$$
zato \v sto je $8 + 4 + 1 = 13$, \v sto zna\v ci da ``3 pi\v sem, a 1 prenosim''.

Pogledajmo sada jedno izvo\dj enje u formalnoj teoriji $\Sab$ da bismo se uverili da ovaj niz pravila
zaista ume da sabere dva broja.

\begin{center}
  \input ft01-sab.pgf
\end{center}



\endinput
%%%%%%%%%%%%%%%%%%%%%%%%%%%%%%%%%%%%%%%%%%%%%%%%%%%%%%%%%%%%%%%%%%%%%%%%%%%%%%%%%%%%%%%

\paragraph{Primer: Formalna teorija BinSab.}
Neka je $\BinSab = (A, F, X, R)$ formalna teorija nad alfabetom $A = \{0, 1, \Boxed\oplus\}$ \v ciji skup formula
je $F = \{\lambda \in A^* : |\lambda|_\oplus \le 1\}$, \v ciji skup aksioma je $X = \{0,1\}^+$ (dakle, sve neprazne re\v ci
nad alfabetom $\{0, 1\}$), i koja ima slede\'cih deset pravila izvo\dj enja, gde je $\lambda, \mu \in X$, a operacija
$\star$ \'ce biti opisana kasnije:
$$
  d : \begin{array}{c}
    \oplus \lambda\\ \hline
    \lambda
  \end{array},\qquad
  p : \begin{array}{c}
    \lambda, \; \mu \; (\text{samo kada je } |\lambda| = |\mu|) \\ \hline
    00 (\lambda \star \mu) 0 \oplus
  \end{array},
$$
$$
  r_0 : \begin{array}{c}
    \lambda 000 \oplus \mu\\ \hline
    \lambda 0 \oplus 0 \mu
  \end{array}, \qquad
  r_1 : \begin{array}{c}
    \lambda 001 \oplus \mu\\ \hline
    \lambda 0 \oplus 1 \mu
  \end{array}, \qquad
  r_2 : \begin{array}{c}
    \lambda 010 \oplus \mu\\ \hline
    \lambda 0 \oplus 1 \mu
  \end{array}, \qquad
  r_3 : \begin{array}{c}
    \lambda 011 \oplus \mu\\ \hline
    \lambda 1 \oplus 0 \mu
  \end{array},
$$
$$
  r_4 : \begin{array}{c}
    \lambda 100 \oplus \mu\\ \hline
    \lambda 0 \oplus 1 \mu
  \end{array}, \qquad
  r_5 : \begin{array}{c}
    \lambda 101 \oplus \mu\\ \hline
    \lambda 1 \oplus 0 \mu
  \end{array},\qquad
  r_6 : \begin{array}{c}
    \lambda 110 \oplus \mu\\ \hline
    \lambda 1 \oplus 0 \mu
  \end{array}, \qquad
  r_7 : \begin{array}{c}
    \lambda 111 \oplus \mu\\ \hline
    \lambda 1 \oplus 1 \mu
  \end{array}.
$$
Operacija $\star$ je definisana na re\v cima $\alpha = a_1 a_2 \ldots a_k$ i $\beta = b_1 b_2 \ldots b_k$
iste du\v zine na slede\'ci na\v cin:
$
  \alpha \star \beta = a_1 b_1 a_2 b_2 \ldots a_k b_k
$.


